hyperparameter selection. All protocol-matched methods share the same
candidate budgets and checkpoint-selection rule. The task-specific
search spaces and selected values are reported in
Appendices~\ref{app:d4rl_methods} and
\ref{app:countdown_methods}.


\subsection{D4RL Baselines and Experimental Protocol}
\label{app:d4rl_methods}

The definitions of Positive-only, Global-\(\alpha\),
Reciprocal-Linear, Reciprocal-Quadratic, and DRPO follow
Appendix~\ref{app:shared_methods}. This section specifies only the
D4RL-specific baselines, remoteness construction, training protocol,
and evaluation procedure.

\paragraph{Tasks and datasets.}

We evaluate HalfCheetah, Hopper, and Walker2d on the medium,
medium-replay, and medium-expert datasets, producing the nine tasks
reported in Table~\ref{tab:d4rl_performance}. The same fixed offline
dataset is used by every protocol-matched method within a task.

\paragraph{Published offline-RL baselines.}

CQL, TD3+BC, and IQL are included as published reference results.
These values are not protocol-matched reruns and are therefore marked
by \(\dagger\) in Table~\ref{tab:d4rl_performance}. The main table
reports their published task-level means. Their original uncertainty
statistics and exact score provenance are recorded separately from the
protocol-matched DRPO experiments.

\paragraph{Gaussian-policy remoteness.}

For an action \(a\in\mathbb{R}^{d_a}\) and current Gaussian policy with
mean \(\mu_\theta(s)\) and standard deviation \(\sigma_\theta(s)\), we use the
following implementation-level standardized remoteness statistic

\[
D_{\mathrm{D4RL}}(s,a)
=
\frac{1}{d_a}
\sum_{j=1}^{d_a}
\left(
\frac{a_j-\mu_{\theta,j}(s)}
     {\sigma_{\theta,j}(s)}
\right)^2.
\]

This dimension-normalized convention differs from
Equation~\eqref{eq:hopper_taper_remoteness} only by a fixed task-level scale.
The associated threshold and scale parameters absorb that constant, and each
protocol keeps its convention fixed across all compared methods.

When the actor uses a transformed action distribution, this quantity is
computed in the policy's native Gaussian coordinates. The resulting
remoteness is used only to determine the negative-update coefficient;
the distance term is detached during the actor update.

\paragraph{Protocol-matched actor comparison.}

Positive-only, Global-\(\alpha\), Reciprocal-Linear,
Reciprocal-Quadratic, and DRPO share the same offline dataset, critic
checkpoint or critic-training protocol, actor initialization, network
architecture, optimizer, minibatch schedule, actor-update budget,
evaluation schedule, and random seeds. They differ only in the
negative-update weighting rule defined in
Appendix~\ref{app:shared_methods}.

The common actor advantage construction is held fixed across these
methods. No method is allowed to change the critic, relabel samples, or
alter the data distribution in order to improve its actor result.

\paragraph{Hyperparameters.}

The registered training and selection configuration is summarized in
Table~\ref{tab:d4rl_registered_config}. All entries must be populated
directly from the frozen experiment configuration rather than
reconstructed after observing the test results.

\begin{table}[t]
\centering
\caption{Registered configuration for the D4RL performance experiments.}
\label{tab:d4rl_registered_config}

\small
\setlength{\tabcolsep}{8pt}
\renewcommand{\arraystretch}{1.12}

\begin{tabular}{@{}ll@{}}
\toprule
Configuration & Registered value \\
\midrule
Actor-update budget
& \textbf{[registered value]} \\
Evaluation interval
& \textbf{[registered value]} \\
Evaluation episodes
& \textbf{[registered value]} \\
Random seeds
& \textbf{[registered seed set]} \\
Near-field threshold candidates, $\tau$
& \textbf{[registered grid]} \\
Remoteness-scale candidates, $c$
& \textbf{[registered grid]} \\
Taper-strength candidates, $\lambda$
& \textbf{[registered grid]} \\
Global-scaling candidates, $\alpha$
& \textbf{[registered grid]} \\
\bottomrule
\end{tabular}
\end{table}

\paragraph{Evaluation and uncertainty.}

Task performance is reported using the standard D4RL normalized return.
For every protocol-matched method, the main table reports the mean over
the registered seeds, while seed-level values and uncertainty
statistics are reported in the appendix. The D4RL-9 total is the sum of
the nine task-level means.

Task-performance collapse is assessed from normalized return and
training trajectories. Variance or support-boundary events are reported
separately from task performance. NaN/Inf numerical failures are
reported as a third outcome and are not merged with either low return or
policy-boundary events. Any claim concerning convergence, collapse, or
method ranking is based on the registered terminal audit rather than an
intermediate checkpoint.


\subsection{Countdown Baselines and Experimental Protocol}
\label{app:countdown_methods}

Positive-only, Global-\(\alpha\), Reciprocal-Linear,
Reciprocal-Quadratic, and DRPO follow the shared definitions in
Appendix~\ref{app:shared_methods}. This section defines the two
published off-policy baselines and the Countdown-specific response-bank,
remoteness, training, and evaluation protocols.

\paragraph{AsymRE.}

Asymmetric REINFORCE defines the response-level advantage using a
tunable reward baseline \(V\):

\[
A_{\mathrm{AsymRE}}(x,y)
=
r(x,y)-V,
\]

and optimizes

\[
\mathcal{J}_{\mathrm{AsymRE}}(\theta)
=
\mathbb{E}_{(x,y)\sim\mathcal{B}}
\left[
\bigl(r(x,y)-V\bigr)
\log\pi_\theta(y\mid x)
\right].
\]

For a binary verifier reward, lowering \(V\) emphasizes successful
responses, whereas increasing \(V\) strengthens suppression of failed
responses. The baseline is selected using the validation protocol and
is held fixed within each run.

Global-\(\alpha\) and AsymRE remain separate comparisons because they
operate on different underlying learning signals. Global-\(\alpha\)
scales the negative part of the common precomputed signed advantage,
whereas AsymRE constructs its signed advantage directly from the raw
verifier reward and the published baseline parameterization. Under
binary rewards and simplified normalization they are closely related,
but the formal comparison preserves their native objectives and tuning
rules.

\paragraph{TOPR.}

Let \(\mu(y\mid x)\) denote the frozen response-generating policy and

\[
\rho_\theta(x,y)
=
\frac{\pi_\theta(y\mid x)}
     {\mu(y\mid x)}
\]

the sequence-level importance ratio. Let
\(R(x,y)\) denote the signed reward used by TOPR and define

\[
\mathcal{T}^{+}
=
\{(x,y):R(x,y)\geq 0\},
\qquad
\mathcal{T}^{-}
=
\{(x,y):R(x,y)<0\}.
\]

We implement canonical TOPR, whose expected update is

\[
\begin{aligned}
g_{\mathrm{TOPR}}(\theta)
={}&
\mathbb{E}_{(x,y)\sim\mu}
\Big[
\mathbf{1}\{(x,y)\in\mathcal{T}^{+}\}
R(x,y)
\nabla_\theta\log\pi_\theta(y\mid x)
\Big]
\\
&+
\mathbb{E}_{(x,y)\sim\mu}
\Big[
\mathbf{1}\{(x,y)\in\mathcal{T}^{-}\}
\min\{\rho_\theta(x,y),1\}
R(x,y)
\nabla_\theta\log\pi_\theta(y\mid x)
\Big].
\end{aligned}
\]

Thus, positive responses receive the SFT-style unit coefficient, while
negative responses are attenuated using a truncated importance ratio.
The behavior-policy log probability required by the ratio is stored
when the frozen response bank is constructed.

\paragraph{Countdown remoteness.}

For a prompt \(x\) and completion
\(y=(y_1,\ldots,y_T)\), learner-relative remoteness is the current mean
completion-token surprisal

\[
D_{\mathrm{CD}}(x,y)
=
-\frac{1}{T}
\sum_{t=1}^{T}
\log
\pi_\theta
\left(
y_t\mid x,y_{<t}
\right).
\]

Only completion tokens are included; prompt tokens and padding positions
are excluded. Mean rather than summed surprisal is used so that response
length does not mechanically determine remoteness. The remoteness score
is detached before applying the shared taper functions.

\paragraph{Frozen response-bank construction.}

The response bank is generated before policy optimization and then held
fixed. Each record contains the puzzle, generated response, verifier
outcome, response length, and the behavior-policy log probability
required by methods that use importance ratios. The same bank and
training split are supplied to every protocol-matched method.

A response is marked successful only when it forms a valid arithmetic
expression, uses the supplied numbers according to the task rules, and
evaluates to the target. Invalid syntax, illegal number use, and
incorrect arithmetic results remain distinct verifier outcomes in the
stored diagnostics, even when they share the same binary training
reward.

\paragraph{Training and checkpoint selection.}

All methods start from the same SFT checkpoint and share the optimizer,
learning-rate schedule, minibatch construction, number of optimizer
steps, validation split, test puzzles, and random seeds. Hyperparameters
are selected using the registered validation rule. The held-out test set
is evaluated only after checkpoint selection.

The registered training and selection configuration is summarized in
Table~\ref{tab:countdown_registered_config}. All entries must be
populated directly from the frozen experiment configuration rather than
reconstructed after observing the test results.

\begin{table}[t]
\centering
\caption{Registered configuration for the Countdown performance experiments.}
\label{tab:countdown_registered_config}

\small
\setlength{\tabcolsep}{8pt}
\renewcommand{\arraystretch}{1.12}

\begin{tabular}{@{}ll@{}}
\toprule
Configuration & Registered value \\
\midrule
Optimizer-step budget
& \textbf{[registered value]} \\
Batch size
& \textbf{[registered value]} \\
Learning-rate schedule
& \textbf{[registered value]} \\
Random seeds
& \textbf{[registered seed set]} \\
AsymRE baseline candidates, $V$
& \textbf{[registered grid]} \\
Near-field threshold candidates, $\tau$
& \textbf{[registered grid]} \\
Remoteness-scale candidates, $c$
& \textbf{[registered grid]} \\
Taper-strength candidates, $\lambda$
& \textbf{[registered grid]} \\
Global-scaling candidates, $\alpha$
& \textbf{[registered grid]} \\
Number of samples for pass@$k$
& \textbf{[registered value]} \\
\bottomrule
\end{tabular}
\end{table}

\paragraph{Verifier-based evaluation.}

Greedy verifier success is the fraction of held-out puzzles solved by
greedy decoding. Pass@\(k\) is the fraction solved by at least one of
the registered \(k\) sampled responses. Valid-expression rate is the
fraction of generated responses that satisfy the parser and number-use
constraints, irrespective of whether they reach the target.

These three metrics are reported separately. In particular, a higher
verifier success rate is not attributed to improved reasoning when it
is accompanied by a deterioration in expression validity.

Task-performance degradation, token- or sequence-support boundary
events, and NaN/Inf numerical failures are audited independently.
Claims about stability or final method ranking use the registered
terminal checkpoint audit rather than finite-step pilot behavior.

\end{document}